The project shall be developed and maintained according to the following general principles.

\subsection{Correctness}

Implementation decisions shall prioritize satisfying the documented requirements and preserving correct behavior. Code and documentation shall not rely on assumptions that are inconsistent with the project requirements.

\subsection{Clarity}

Code, configuration, and documentation shall be written so that they can be understood by an appropriately skilled contributor who was not the original author. Clear names, simple structures, consistent formatting, and concise documentation should be preferred over clever or unnecessarily complex solutions.

\subsection{Consistency}

Similar problems shall be solved in similar ways unless a documented reason justifies a different approach. Project conventions shall be applied consistently across source code, documentation, directory structures, configuration, and tooling.

\subsection{Maintainability}

Designs should minimize unnecessary coupling and complexity. Components should have clear responsibilities, stable interfaces, and appropriate separation of concerns. Changes should be made in a way that reduces the long-term cost of maintenance.

\subsection{Testability}

Important behavior shall be designed so that it can be verified. New or modified functionality should include appropriate tests, checks, or other evidence that the implementation satisfies its requirements.

\subsection{Security}

Security-relevant requirements shall be considered throughout the development lifecycle. Sensitive data shall be protected, access shall be limited to what is necessary, and unsafe defaults shall be avoided.

\subsection{Portability and Reproducibility}

The project should avoid unnecessary dependence on a particular computer, operating system, compiler, interpreter, tool version, or local configuration. Builds, tests, and other important activities should be reproducible using documented dependencies and procedures.

\subsection{Traceability}

Requirements, implementation decisions, tests, defects, reviews, and releases should be traceable where practical. The level of traceability shall be appropriate to the size, risk, and purpose of the project.

\subsection{Automation}

Repetitive checks should be automated where practical. Formatting checks, static analysis, tests, documentation generation, and other suitable verification activities should be integrated into the normal development workflow.

\subsection{Small and Reviewable Changes}

Changes should be limited to a clear purpose and kept small enough to review effectively. Unrelated refactoring, formatting changes, or file movements should not be combined with functional changes unless there is a clear benefit.

\subsection{Documentation as a Project Artifact}

Documentation shall be treated as part of the project rather than as an optional supplement. Documentation affected by a change shall be updated as part of that change.